\subsubsection*{引理 3.6 (Euler).}
\underline{设 $(x_0,H_0)$ 为 $\mathbb P(N)$ 的闭点, 以 $T_L$ 表示由 $x_0$ 定义的超平面 $L\subset\check{\mathbb P}^{r}$ 的切丛. 则在 $T_{\check{\mathbb P}^{r},H_0}$ 中有}
\[
d\varphi\bigl(T_{\mathbb P(N),(x_0,H_0)}\bigr)
\subset T_{L,H_0}.
\]

\noindent\underline{证明.}
选取 $\mathbb P^r$ 上的齐次坐标 $X_0,\ldots,X_r$, 使 $x_0$ 为点 $\underline A=(A_0,\ldots,A_r)$. $T_{X,x_0}$ 的一个点是 $\mathbb P^r$ 的一个取值于 $k[\epsilon]/(\epsilon^2)$ 的点, 形如
\begin{equation}\tag{3.6.1}
\underline A+\epsilon\underline B,
\end{equation}
且对 $I$ 的每个局部截面 $H$, 有
\begin{equation}\tag{3.6.2}
\sum_{i=0}^{r}B_i\frac{\partial H}{\partial X_i}(\underline A)=0.
\end{equation}
再选取 $I$ 的一个局部截面 $H$, 使它经 (3.1.1) 给出 $H_0$. 由于 $\mathbb P(N)$ 在 $X$ 上局部平凡, $T_{\mathbb P(N),(x_0,H_0)}$ 的一个点可表示为
\begin{equation}\tag{3.6.3}
(\underline A+\epsilon\underline B,\ H+\epsilon F),
\end{equation}
其中 $F$ 为 $I$ 的局部截面, $\underline A,\underline B$ 如上. 它在 $\varphi$ 下的像为 $k[\epsilon]/(\epsilon^2)$ 上由方程
\begin{equation}\tag{3.6.4}
\sum_{i=0}^{r}
\frac{\partial(H+\epsilon F)}{\partial X_i}
(\underline A+\epsilon\underline B)T_i=0.
\end{equation}
定义的超平面. 因而所需断言等价于
\begin{equation}\tag{3.6.5}
\sum_{i=0}^{r}A_i
\frac{\partial(H+\epsilon F)}{\partial X_i}
(\underline A+\epsilon\underline B)=0.
\end{equation}
由于 $\epsilon^2=0$, Taylor 公式给出
\begin{equation}\tag{3.6.6}
\frac{\partial(H+\epsilon F)}{\partial X_i}
(\underline A+\epsilon\underline B)
=\frac{\partial H}{\partial X_i}(\underline A)
+\epsilon\frac{\partial F}{\partial X_i}(\underline A)
+\epsilon\sum_{j=0}^{r}B_j
\frac{\partial^2H}{\partial X_j\partial X_i}(\underline A).
\end{equation}
因此只需检验
\begin{equation}\tag{3.6.7}
\sum_{i=0}^{r}A_i\frac{\partial H}{\partial X_i}(\underline A)=0,
\end{equation}
\begin{equation}\tag{3.6.8}
\sum_{i=0}^{r}A_i\frac{\partial F}{\partial X_i}(\underline A)=0,
\end{equation}
以及
\begin{equation}\tag{3.6.9}
\sum_{i,j=0}^{r}A_iB_j
\frac{\partial^2H}{\partial X_i\partial X_j}(\underline A)=0.
\end{equation}
此处 $H$ 和 $F$ 由零次齐次函数表示, 因此 Euler 公式给出
\begin{equation}\tag{3.6.10}
\sum_iX_i\frac{\partial H}{\partial X_i}=0,
\qquad
\sum_iX_i\frac{\partial F}{\partial X_i}=0.
\end{equation}
由于 $\partial H/\partial X_j$ 为 $-1$ 次齐次函数, 同一公式给出
\[
\sum_iX_i\frac{\partial^2H}{\partial X_i\partial X_j}
=-\frac{\partial H}{\partial X_j}.
\]
从而 (3.6.9) 化为
\begin{equation}\tag{3.6.11}
-\sum_jB_j\frac{\partial H}{\partial X_j}(\underline A)=0,
\end{equation}
这恰好就是 (3.6.2).

这就证明 $\varphi$ 诱导同构
\begin{equation}\tag{3.6.12}
\varphi^{-1}(V')\xrightarrow{\sim}V',
\end{equation}
其逆为 $\psi$.